\section*{الفصل \ref{ch:b2:speciation} --- الانتواع والتطور الكبروي}

\begin{solution}{exo:b2:speciation:1}
النوع مجموعة جمهرات طبيعية تتزاوج فيما بينها
وتكون معزولة تكاثريًا عن المجموعات الأخرى من هذا القبيل. وهو يخفق في
السلالات اللاجنسية (البكتيريا، والدولابيات البدلوئية، والهندباء:
فيُستعمل المفهوم الشكلي أو النشوئي، أي عناقيد الأشكال المتشابهة أو
السلالات القابلة للتشخيص)؛ وفي الأحافير (فالمفهوم الشكلي، إذ
لا يمكن إجراء أي تهجين)؛ وفي الجمهرات المتباعدة مكانيًا التي لا
تلتقي قط (فالمفهوم البيئي أو النشوئي، إذ لا يمكن اختبار
التزاوج) --- وهو مشوَّش بالبلوط المتهجن وبالأنواع
الحلقية.
\end{solution}

\begin{solution}{exo:b2:speciation:2}
الضفدعان في شهرين مختلفين: قبل اللاقحة (زماني). والبغل: بعد اللاقحة
(عقم الهجين). والنطفة التي لا تستطيع الارتباط: قبل اللاقحة (عرسي).
والهجائن العافية ذات البذور العقيمة: بعد اللاقحة (عقم الهجين،
أو انهيار الهجين إن كان الإخفاق في الجيل التالي).
وصرصارا الليل بأغنيتين مختلفتين: قبل اللاقحة (سلوكي).
\end{solution}

\begin{solution}{exo:b2:speciation:3}
الجغرافي: حاجز يقسم مدى (جمبري الطقطقة على جهتي
برزخ بنما). والمحيطي: بضعة مؤسسين وراء
المدى (عصافير داروين من سلف في بر رئيس). والمتشارك المكان:
بلا حاجز (ذبابة يرقة التفاح على الزعرور والتفاح؛ وقمح
متعدد الصيغ الغيري). والمتجاور: على امتداد تدرج مع \omterm{prop:b2:speciation:reinforcement}{منطقة
تهجين} في الوسط (الأعشاب على مخلفات المناجم وخارجها، وهي تزهر في
أوقات مختلفة).
\end{solution}

\begin{solution}{exo:b2:speciation:4}
يفوز النوع 1 دائمًا؛ أو يفوز النوع 2 دائمًا؛ أو تعايش مستقر؛
أو يفوز أيهما بحسب الأعداد الابتدائية. والتعايش مستقر
حين يستطيع كل نوع غزو الآخر عند حمولته القصوى،
أي حين يحد كل نوع نفسه أكثر مما يحد
الآخر ($K_1 > \alpha K_2$ و$K_2 > \beta K_1$).
\end{solution}

\begin{solution}{exo:b2:speciation:5}
$\alpha K_2 = 400 < K_1 = 1000$: فالنوع 1 يغزو؛ و$\beta K_1 = 600
< K_2 = 800$: فالنوع 2 يغزو. و$N_1^{*} = (1000 - 400)/(1 - 0.3) =
857$، و$N_2^{*} = (800 - 600)/0.7 = 286$؛ وعند التوازن يكون النوع
1 دون حمولته $K$ بمقدار $\alpha N_2^{*} = 143$ والنوع 2 بمقدار $\beta
N_1^{*} = 514$.
\end{solution}

\begin{solution}{exo:b2:speciation:6}
$\alpha K_2 = 1200 > K_1$: فالنوع 1 لا يستطيع غزو النوع 2 عند
حمولته القصوى؛ و$\beta K_1 = 300 < K_2$: فالنوع 2 يستطيع غزو
النوع 1. فيفوز النوع 2 دائمًا، لأن أفراده يضرون
النوع 1 أكثر مما يضره أفراده هو في حين لا يكاد هو يتأثر.
وعند $\alpha = \beta = 1.5$: يكون $\alpha K_2 = 1200 > 1000$
و$\beta K_1 = 1500 > 800$: فلا يستطيع أي منهما غزو الآخر، ويقصي كل
منهما الآخر إذا استقر، ويفوز المؤسس.
\end{solution}

\begin{solution}{exo:b2:speciation:7}
رباعي الصيغة $4n$ يصنع أعراسًا $2n$؛ وإذا هُجّن بأعراس \omterm{def:b2:meiosis-heredity:meiosis}{ثنائي الصيغة}
$n$ أعطى ثلاثيات الصيغة $3n$، ولا يستطيع تقسمها المنصّف مزاوجة ثلاث
مجموعات فيعطي أعراسًا مختلة الصيغة غير قابلة للحياة في معظمها. و
رباعي الصيغة لا يستطيع التكاثر إلا مع نفسه (أو بالإلقاح الذاتي): فهو معزول في
خطوة واحدة. وهو يخفق عادةً لأنه وحيد --- فشركاؤه المحتملون
الوحيدون ثنائيات الصيغة وكل تهجين كهذا مهدور ---
فيُغمر؛ وهو يدوم إذا ألقح نفسه، أو تكاثر خضريًا،
أو نشأ مرارًا.
\end{solution}

\begin{solution}{exo:b2:speciation:8}
$0.01\times 10\times 10 = 1$ حالة عدم توافق؛ وعند 30 لكل منهما، $0.01\times
900 = 9$. فالعزل ينمو كمربع زمن التباعد: فجمهرتان
تباعدتا ثلاثة أضعاف المدة تكونان أكثر عزلًا بتسعة
أضعاف --- وهي كرة الثلج.
\end{solution}

\begin{solution}{exo:b2:speciation:9}
$w = 2\sqrt{8/0.05} = \qty{25}{km}$؛ و$w = 0.3\sqrt{8/0.4} =
\qty{1.3}{km}$. والتعزيز أرجح في الثانية: فالهجائن
غير لائقة بشدة، فالأنثى التي ترفض الصنف الآخر تكسب
كثيرًا، والمنطقة ضيقة بما يكفي لأن يلتقي الشكلان مرارًا
تكفي لانتقاء ذلك الرفض.
\end{solution}

\begin{solution}{exo:b2:speciation:10}
$\alpha K_2 = 450 < 500$، فيغزو كل منهما الآخر: فيتعايشان،
عند $N^{*} = (500 - 450)/(1 - 0.81) = 263$ لكل منهما --- بالكاد، قرب
حافة الإقصاء، ولو تغيرت $K$ قليلًا لمالت الكفة. وعند
$\alpha = \beta = 0.4$: $N^{*} = (500 - 200)/(1 - 0.16) = 357$ لكل منهما.
فالإزاحة أدارت كل خط تساوٍ بعيدًا عن الآخر (إذ تحرك $K/\alpha$
إلى الخارج)، فابتعد التقاطع عن المحورين وصار كل
نوع أقرب إلى حمولته $K$: أي تعايشًا أوثق مع
أعداد أكبر من كل منهما.
\end{solution}

\begin{solution}{exo:b2:speciation:11}
$R = h^{2}S = 0.7\times(-1) = \qty{-0.7}{mm}$: كما لوحظ. وفي
جفاف سنة 1977، ولم يكن العصفور الأرضي الكبير حاضرًا، انتُقيت العصافير
المتوسطة نفسها لمناقير \emph{أكبر}، إذ كانت البذور الكبيرة
هي الغذاء الباقي؛ وفي سنة 2004 أخذ الغازي البذور الكبيرة،
فخسرت العصافير المتوسطة التي نافست عليها: فاتجاه
الانتقاء قرره المنافس.
\end{solution}

\begin{solution}{exo:b2:speciation:12}
تُصنع بالمصادفة: فنموذج الموقعين يبيّن العزل ناشئًا ناتجًا
عرضيًا لاستبدالات ثُبّتت لأسباب أخرى، أو بالانجراف،
والجغرافيا --- حاجز أو تأسيس --- تمد الفصل
الذي يتيح للمصادفة أن تتراكم. وتُحفظ بالانتقاء: فحيث تلتقي
الأشكال، يبني الانتقاء ضد الهجائن غير اللائقة حواجز قبل
اللاقحة (وهو التعزيز)، وتزيح المنافسة الصفات حتى
يستطيع الاثنان التعايش. والقول يغفل \omterm{prop:b2:speciation:modes}{الانتواع المتشارك المكان}
بالانتقاء المفرِّق، حيث يصنع الانتقاء النوعين كما يحفظهما،
ويغفل تعدد الصيغ، حيث تكون المصادفة هي القصة كلها.
\end{solution}

\begin{solution}{pb:b2:speciation:1}
\textbf{1.} النوع 1: $N_1 + 0.9N_2 = 1200$، وتقاطعاه $N_1 = 1200$
و$N_2 = 1333$. والنوع 2: $N_2 + 0.3N_1 = 400$، وتقاطعاه $N_2 =
400$ و$N_1 = 1333$.
\textbf{2.} $\beta K_1 = 360 < K_2 = 400$: فالنوع 2 يغزو؛
و$\alpha K_2 = 360 < K_1 = 1200$: فالنوع 1 يغزو. أي تعايش
مستقر.
\textbf{3.} $N_1^{*} = (1200 - 360)/(1 - 0.27) = 1151$؛ و$N_2^{*} =
(400 - 360)/0.73 = 55$.
\textbf{4.} $\mathrm{d}N_1/\mathrm{d}t = 0.5\times 1200\,(1 - 1218/1200)
= -9$ في السنة: فالمقيم، عند حمولته القصوى، يدفعه الوافدون
إلى الهبوط.
\textbf{5.} $K_1 = 600$، و$K_2 = 200$: $N_1^{*} = (600 - 180)/0.73 =
575$، و$N_2^{*} = (200 - 180)/0.73 = 27$: فينتصف الاثنان، ويكون الغازي،
عند 27 طائرًا، قريبًا من الانقراض --- فالجفاف هو حين
تعض المنافسة.
\textbf{6.} الغازي هو الطائر الأكبر ذو المنقار الأكبر: فهو
يأخذ البذور الكبيرة التي يستعملها المقيم أيضًا، ويأخذها
أسرع، في حين أن بذور المقيم الصغيرة قليلة النفع له.
\textbf{7.} $S = 9.0 - 9.5 = \qty{-0.5}{mm}$؛ و$R = 0.7\times(-0.5) =
\qty{-0.35}{mm}$.
\textbf{8.} $9.5 - 0.35 = \qty{9.15}{mm}$.
\textbf{9.} $N_1^{*} = (1200 - 200)/(1 - 0.15) = 1176$؛ و$N_2^{*} =
(400 - 360)/0.85 = 47$.
\textbf{10.} ارتفع تقاطع $N_2$ لخط تساوي النوع 1 من
$K_1/0.9 = 1333$ إلى $K_1/0.5 = 2400$: فدار الخط بعيدًا عن
خط الغازي، وصار المقيم أقل تأثرًا بكل غازٍ،
وتحرك التوازن نحو حمولة المقيم $K$.
\textbf{11.} $2/0.35 \approx 6$ سنوات. ولا يحدث ذلك لأن انتقاءً
من هذا القبيل لا يقع إلا في سنوات الجفاف، وينقلب في السنوات الرطبة
حين تكثر البذور الصغيرة (أو حين يندر الغازي)،
ولأن التباين الجمعي سيُستنفد.
\textbf{12.} مليونا سنة نحو $10^{6}$ جيل؛ وإزاحة
قدرها \qty{0.35}{mm} في الجيل لو دامت جزءًا من ألف
منها لغيّرت المنقار بمئات المليمترات. فالانتقاء
قوي لكنه غير ثابت؛ والتغير الموجّه المستدام هو الاستثناء
النادر، والتغير الصافي على المدى الجيولوجي بقية ضئيلة
من حلقات كبيرة متعاكسة.
\textbf{13.} الأغنية، وهي تُتعلَّم من الأب وتستعملها الإناث في
اختيار القرين، مع شكل المنقار نفسه إشارةً. والهجائن ذات
المناقير الوسطية لا تعالج البذور الكبيرة ولا الصغيرة
أحسن معالجة فتجوع أولًا في الجفاف.
\textbf{14.} $w = 0.5\sqrt{8/0.2} = \qty{3.2}{km}$.
\textbf{15.} الجزيرة أضيق مما ستكون عليه المنطقة: فلا
يتشكل أي تدرج مكاني؛ والتهجين، حيث يقع، يمس
الجمهرة كلها، ولا يبقي الاثنين متفرقين إلا التزاوج
الانتقائي.
\textbf{16.} الإناث التي لا تتزاوج إلا مع ذكور من أغنيتها
ومنقارها تترك أحفادًا أكثر ممن تتهجن؛ فيصير التفضيل
والإشارة (الأغنية والمنقار) أشد تميزًا حيث يعيش
النوعان معًا منه حيث يعيشان متفرقين.
\textbf{17.} $0.005\times 15\times 15 = 1.1$ حالة عدم توافق: أي
واحدة وسطيًا، مع هجائن مرجح أنها منخفضة \omterm{def:b2:population-genetics:fitness}{اللياقة} لا عقيمة. فليسا
نوعين بعدُ بالمعيار الحيوي؛ لكنهما في الطريق.
\textbf{18.} الحيوي: ليس بعد (فالهجائن شبه لائقة).
والشكلي: ربما، إن كانت المناقير أو الريش قد تباعدت
بصورة مرئية. والنشوئي: نعم، إن كان لكل جمهرة فرق
تشخيصي مثبَّت. والمفاهيم تختلف بالضبط في أثناء
عملية الانتواع، وهو ما يجعلها مفيدة.
\textbf{19.} $1/(5\times 10^{6}) = 2\times 10^{-7}$ للنوع في
السنة؛ ولأجل $10^{7}$ نوع، انقراضان في السنة.
\textbf{20.} من 200 إلى 2000 نوع في السنة.
\textbf{21.} $0.75\times 10^{7}/10^{5} = 75$ نوعًا في السنة --- أي أقل
من المعدل الحالي عند حده الأدنى. فالحلقة الحالية، إن
دامت، \omterm{prop:b2:speciation:tempo}{انقراض جماعي} يمضي أسرع من الخمسة
الكبار.
\textbf{22.} المكانات التي أخلاها الخاسرون خالية من
المنافسين، فيُحبَّذ أي متغير ناجٍ يستطيع استعمال إحداها،
فتتنوع الناجيات فيها.
\textbf{23.} \omterm{prop:b2:speciation:tempo}{التوازن المتقطع}؛ \omterm{prop:b2:speciation:modes}{والانتواع الجغرافي} (المحيطي)
في جمهرة صغيرة معزولة ينتشر سليلها بعد ذلك
في مدى الجمهرة الرئيسة.
\textbf{24.} $1.04^{n} = 2$: أي $n = \ln 2/\ln 1.04 = 18$ جيلًا.
والسجل يُظهر ملايين السنين لأن الانتقاء ينقلب،
وينهض \omterm{prop:b2:population-genetics:kinds}{الانتقاء المثبِّت} بمعظم الوقت، والتغير الصافي
هو الفرق الصغير بين حلقات كبيرة متعاكسة.
\textbf{25.} قبل الإزاحة: تعايش مستقر عند $N_1^{*} =
1151$ و$N_2^{*} = 55$؛ وبعدها: 1176 و47. وإزاحة المنقار
\qty{-0.35}{mm} في جيل واحد. وعرض \omterm{prop:b2:speciation:reinforcement}{منطقة التهجين} \qty{3.2}{km}.
\end{solution}
